\documentclass{article}
\usepackage{amsmath}
\usepackage{amsthm}
\usepackage{amsfonts}
\usepackage{amssymb}
\usepackage{bbold}
\usepackage{dsfont}
\usepackage{graphicx}
\usepackage{cleveref}
\usepackage{bm}
\usepackage{booktabs}
\usepackage{siunitx}
\usepackage{xcolor}
\usepackage[utf8]{inputenc}
\usepackage[english]{babel}
\newtheorem{definition}{Definition}
\newtheorem{example}{Example}
\newcommand{\exP}[1]{\mathrm{e}^{#1}}
\newcommand{\dif}{\mathrm{d}}
\renewcommand{\vec}[1]{\mathbf{#1}}
\newcommand{\pdagger}{{\phantom{\dagger}}}
\newcommand{\state}[1]{\vert #1\rangle}
\newcommand{\cstate}[1]{\langle #1\vert}
\newcommand{\tr}[1]{\mathrm{tr}\left\{#1\right\}}
\newcommand{\trs}[2]{\, \mathrm{tr}_{#1}\{#2\}}
\newcommand{\expc}[3]{\langle #1\vert #2 \vert #3\rangle}
\usepackage[doi=false,url=false,isbn=false,backend=bibtex,style=nature,natbib=true,sorting=none]{biblatex}
\bibliography{bib.bib}
\usepackage{geometry}
\geometry{
	a4paper,
	total={170mm,257mm},
	left=20mm,
	top=20mm,
}
\begin{document}
\title{Analysis of steady states within Ehrenfest-HEOM}
\author{Rudolf Smorka}
\maketitle
\section{Aim}
The aim of this article is to analyze the steady states of the Ehrenfest-HEOM numerically. 
\section{Equations of motion}
The Ehrenfest-HEOM is a nonlinear system which reads 
\begin{align}
	\frac{\dif}{\dif t}\rho(t) &= -i[H_\mathrm{el}+H_\mathrm{el-vib}(X(t)),\rho] + \sum_\ell (\Pi_\ell^\dagger + \Pi_\ell^{\phantom{\dagger}}) \\
	\Pi_\ell(t) &= \frac{1}{4}(\mathds{1}-2\rho)\Gamma_\ell-\sum_{p=1}^{N_p} \Pi_{\ell,p}(t)\\
	\frac{\dif}{\dif t} \Pi_{\ell,p}(t) &= -\frac{i\eta_p}{\beta_\ell}\Gamma_\ell - i\left(H_\mathrm{el}+H_\mathrm{el-vib}(X(t))-\frac{i}{2}\Gamma-\chi_p^+\mathds{1}\right)\Pi_{\ell,p}(t)\\
	\frac{\dif P_j}{\dif t}&= -X_j - \tr{\rho(t) \nabla_{X_j}H_\mathrm{el-vib}}, \\
	\frac{\dif X_j}{\dif t} &=  P_j.
\end{align}
The EOM of the reduced density matrix can be rewritten to
\begin{align}
		\frac{\dif}{\dif t}\rho(t) &= -i[H_\mathrm{el}+H_\mathrm{el-vib}(X(t)),\rho] \\
		&+ \sum_\ell \left[ \frac{1}{4}(\mathds{1}-2\rho)\Gamma_\ell-\sum_{p=1}^{N_p} \Pi_{\ell,p}(t) + \left(\frac{1}{4}(\mathds{1}-2\rho)\Gamma_\ell-\sum_{p=1}^{N_p} \Pi_{\ell,p}(t)\right)^\dagger \right]\\
		&=-i[H_\mathrm{el}+H_\mathrm{el-vib}(X(t)),\rho] \\
		&+ \sum_\ell \left[ \frac{1}{2}[\Gamma_\ell-(\rho\Gamma_\ell+\Gamma_\ell\rho)]-\sum_{p=1}^{N_p} \left(\Pi_{\ell,p}(t) + \Pi^\dagger_{\ell,p}(t)\right)\right]\\
		\overset{\mathrm{vectorization}}{\Rightarrow}  \frac{\dif}{\dif t}\vert \rho(t)\rangle &= (\mathcal{L}_\mathrm{el}+\mathcal{L}_\mathrm{el-vib}(X(t)) - \frac{1}{2}\sum_\ell[ \Gamma_\ell^T\otimes \mathds{1}+\mathds{1}\otimes \Gamma_\ell])\vert \rho\rangle + \sum_\ell\left[ \frac{1}{2}\vert\Gamma_\ell\rangle - \sum_{p=1}^{N_p} \left(\vert \Pi_{\ell,p}(t)\rangle + \vert \Pi^\dagger_{\ell,p}(t)\rangle\right)\right]\\
		&\equiv \mathcal{A}(X(t)) \vert \rho\rangle + \vert \mathcal{B}(t)\rangle,
\end{align}
since $\rho=\rho^\dagger$ and $\Gamma_\ell = \Gamma_\ell^\dagger$, and $\ell = L,R$. Here we defined the Liouvillian $\mathcal{L}_\mathrm{el}\rho = -i[H_\mathrm{el},\rho]$, and for $H_\mathrm{el-vib}$ analogously. Its vectorized form takes $\mathcal{L}\vert \rho\rangle = -i ((\mathds{1}\otimes H)-(H^T\otimes \mathds{1}))\vert \rho\rangle$. Further, $\mathcal{A}(X(t))$ is the superoperator 
\begin{equation}
	\mathcal{A}(X(t)) = \mathcal{L}_\mathrm{el}+\mathcal{L}_\mathrm{el-vib}(X(t)) - \frac{1}{2}\sum_\ell[ \Gamma_\ell^T\otimes \mathds{1}+\mathds{1}\otimes \Gamma_\ell],
\end{equation}
and the vector $\vert \mathcal{B}(t)\rangle$ is defined by
\begin{equation*}
	\vert \mathcal{B}(t)\rangle )=\sum_\ell\left[\frac{1}{2}\vert\Gamma_\ell\rangle -  \sum_{p=1}^{N_p} \left(\vert \Pi_{\ell,p}(t)\rangle + \vert \Pi^\dagger_{\ell,p}(t)\rangle\right)\right].
\end{equation*}
Next, the EOM for the Pad\'{e} resolved auxiliary density matrices is vectorized
\begin{align}
	\frac{\dif}{\dif t}\vert \Pi_{\ell,p}(t)\rangle &= -\frac{i\eta_p}{\beta_\ell}\vert\Gamma_\ell\rangle - i\mathds{1}\otimes \left(H_\mathrm{el}+H_\mathrm{el-vib}(X(t))-\frac{i}{2} \Gamma-\chi_p^+\mathds{1}\right)\vert \Pi_{\ell,p}(t)\rangle\\
	&\equiv -\frac{i\eta_p}{\beta_\ell}\vert\Gamma_\ell\rangle + \mathcal{K}_{p}(X(t)).
\end{align}
Now, the combined vector of density matrices is introduced $\vert \sigma\rangle = \vert \rho\rangle \oplus \bigoplus_{\ell}\bigoplus_{p}\vert \Pi_{\ell,p}\rangle$, such that the electronic part of the equations of motion is reduced to a single matrix equation
%\begin{scriptsize}
\begin{align*}
	\frac{\dif}{\dif t}\vert \sigma(t)\rangle &= \begin{pmatrix}
		\mathcal{A}(X(t)) & -\mathds{1}_{n^2}-\mathds{1}^{\dagger}_{n^2}  & \dots & \dots& -\mathds{1}_{n^2}-\mathds{1}^{\dagger}_{n^2} & -\mathds{1}_{n^2}-\mathds{1}^{\dagger}_{n^2}\\
		\mathbb{0} &  \mathcal{K}_1(X(t))  & \dots &\dots &\mathbb{0}  &\mathbb{0}\\
		\vdots 	& & \ddots &  & &\vdots\\
		 \dots & \mathbb{0} & \mathcal{K}_{N_p}(X(t)) & \mathbb{0} & \dots & \\
		\vdots & \dots & \mathbb{0}  & \mathcal{K}_{1}(X(t)) & \mathbb{0} & \vdots \\
		\vdots 	& & & \ddots &  &\vdots\\		
		\mathbb{0} & \dots & & \dots & \mathbb{0}& \mathcal{K}_{N_p}(X(t))
	\end{pmatrix}\vert \sigma\rangle 
+ \begin{pmatrix}
	\vert\mathds{1}\rangle\\
	-\frac{i\eta_1}{\beta_L}\vert\Gamma_L\rangle\\
	\dots\\
	-\frac{i\eta_{N_p}}{\beta_L}\vert\Gamma_L\rangle\\
	-\frac{i\eta_1}{\beta_R}\vert\Gamma_R\rangle\\
	\dots\\
	-\frac{i\eta_{N_p}}{\beta_R}\vert\Gamma_R\rangle
\end{pmatrix}\\
&\equiv \mathcal{M}(X(t)) \vert \sigma \rangle + b.
\end{align*}
%\end{scriptsize}
where $n$ denotes the linear dimension of the original problem ($n=2L$ with $L$ lattice sites), $\mathds{1}^\dagger \vert A\rangle = \vert A^\dagger \rangle$.
Note the particular structure of $\mathcal{M}(X(t))$: It has non-zero matrices only in the first row and on the diagonal. 

\textit{1) Constant $X(t)$:} Consider first a fixed $X(t)=X=\mathrm{const}$. In this case, one has a linear, unforced, nonhomogeneous system. The system for $\vert \Pi_{\ell,p}\rangle$ has the solution $\vert \Pi_{\ell,p}(t)\rangle$ for the initial condition $\vert \Pi_{\ell,p}(t_0)\rangle$, and this solution defines $\vert \mathcal{B}(t)\rangle$ for all times. The formal solution of the reduced density matrix part is
\begin{equation*}
\vert \rho(t)\rangle = \exP{\mathcal{A}(X)t} \int_0^t \exP{-\mathcal{A}(X)s}\vert \mathcal{B}(s)\rangle \dif s \equiv \exP{\mathcal{A}(X)t} \vert F(t)\rangle.
\end{equation*}
The steady state can be inferred from solving the kernel of the system. From the equation of motion, one obtains the condition
\begin{equation*}
	0 = \mathcal{A}\vert \rho\rangle + \vert \mathcal{B}(t)\rangle 
\end{equation*}
and from the formal solution one gets
\begin{align*}
	0 &=  \mathcal{A}\exP{\mathcal{A}t} \vert F(t)\rangle + \exP{\mathcal{A}t} \frac{\dif }{\dif t}F(t)\\
	&= \mathcal{A}\exP{\mathcal{A}t} \vert F(t)\rangle + \exP{\mathcal{A}t}\exP{-\mathcal{A}t}\vert \mathcal{B}(t)\rangle\\
	&= \mathcal{A}\exP{\mathcal{A}t} \vert F(t)\rangle +\vert \mathcal{B}(t)\rangle.
\end{align*}
Equating both gives 
\begin{equation*}
	\mathcal{A}\vert \rho\rangle = \mathcal{A}\exP{\mathcal{A}t} \vert F(t)\rangle
\end{equation*}
resulting in the steady state density matrix
\begin{equation*}
	\vert \rho_\mathrm{ss} \rangle = \lim_{t\rightarrow \infty} \exP{\mathcal{A}t}\int_0^t \exP{-\mathcal{A}s}\vert \mathcal{B}(s)\rangle \dif s.
\end{equation*}

\printbibliography
\end{document}

